\subsection{Reducción a un permanente auxiliar}\label{costes:permanente}

Sea $\varepsilon_{\rm out}\ge\max_f\ell_f$ la cota exterior común, con $\varepsilon_{\rm out}<1$.
Sea $P=(p_{ij})$ la matriz del punto estacionario interior de Bethe, con ceros fuera del soporte, y sea
$W_{ij}=p_{ij}(1-p_{ij})$. En lo que sigue $W$ denota esta matriz auxiliar. Escribimos
\begin{equation}\label{coste:rook-def}
 \begin{gathered}
 S_{R,d}=A_R^T\operatorname{diag}(\tau_d(p_e))A_R,\qquad S_2=S_{R,2},\\
 C_2=\sum_{ij}\bigl[\log(1-2p_{ij})-2\log(1-p_{ij})\bigr],\qquad
 \Gamma_2=\frac{e^{-C_2/2}\sqrt{4\pi k}}{\sqrt{\pdet S_2}} .
 \end{gathered}
\end{equation}
Una permutación en el soporte usa una casilla por fila y por columna. La llamamos legal si además usa
cada diagonal a lo sumo una vez.

\begin{lema}[Reducción a un permanente]\label{coste-factor}
Sea $\mu_W(\sigma)=\prod_{e\in\sigma}W_e/\per W$ la probabilidad que la matriz $W$ induce en las
permutaciones del soporte, y sea
\[
 F(\sigma)=
 \frac{\prod_f h_f}{\prod_e(1-p_e)^2}\,
 \mathbf1_{\{\sigma\ {\rm legal}\}}
 \prod_{e\in\sigma}\frac{(1-p_e)^2}{h_{D(e)}h_{A(e)}} .
\]
Entonces
\[
 Q(B)e^{-H_B(p)}
   =\frac{\per W}{\prod_e(1-p_e)}\,\E_{\mu_W}F.
\]
Para $Q(B)>0$, supongamos que la matriz bistocástica $P$ cumple las cotas de
\eqref{terminal-hipotesis-matriz} con $K=k$, es decir, $p_{ij}\le M_0/k$ y
$\|P-\mathbf1\mathbf1^T/k\|_{\rm op}\le\rho_*<1$. Entonces, uniformemente para $M_0,\rho_*$ fijos,
\[
 g(B)=\log\E_{\mu_W}F-\log(\Gamma_4/\Gamma_2)+o(1)
 \qquad (k\longrightarrow\infty).
\]
\end{lema}
\begin{proof}
Las filas y columnas de $P$ suman uno. La masa de cualquier conjunto de filas está contenida en sus
columnas vecinas, cada una de capacidad uno. Por el teorema de Hall existe una permutación del soporte
y $\per W>0$. En una permutación legal, multiplicar $W_e$ por la fracción de la última productoria da
\[
 \frac{p_e(1-p_e)^3}{h_{D(e)}h_{A(e)}}=a_i b_j
\]
por el Lema~\ref{coste-vacancias}. El producto sobre la permutación vale entonces $\prod_i a_i\prod_j
b_j$, independientemente de la colocación. Sumando sobre las $Q(B)$ permutaciones legales, el lado
derecho de la primera igualdad es
$Q(B)\mathcal K/\prod_e(1-p_e)^3$, que vale $Q(B)e^{-H_B(p)}$ por ese mismo lema.

Para la fórmula asintótica aplicamos la estimación matricial del Lema~\ref{terminal-normalizador} a
esta $P$, con $K=k$, $W_{ij}=p_{ij}(1-p_{ij})$ y el factor de torres de \eqref{coste:rook-def}. Se
obtiene $\log(\per W/\prod_e(1-p_e))=\log\Gamma_2+o(1)$. Su prueba usa las estimaciones del escalado
y del permanente de la Sección~\ref{permanente}. En la ecuación estacionaria ponderada de torres,
$p_e(1-p_e)=W_e$, todos los multiplicadores pueden tomarse iguales a uno. Tomar logaritmos en la
primera igualdad y usar $L(B)=H_B(p)+\log\Gamma_4$ prueba la
segunda. La ley $\mu_W$ se usa únicamente para contar en $B$.
\end{proof}

\subsection{Los términos del cociente de normalizadores}

\begin{lema}[Inversa de las matrices de filas y columnas]\label{coste-inversa}
Bajo las hipótesis del Lema~\ref{coste-estatico}, escribamos
$S_0=A_R^T\operatorname{diag}(p)A_R$ y $\iota_e$ para el vector indicador de la fila y la columna de
$e$.
Para $0\le s\le1$ ponemos
$S(s)=(1-s)S_{R,2}+sS_{R,4}$. Uniformemente en $s$ y en los índices físicos $u,v$,
\[
 (S(s)^+)_{uv}=\mathbf1_{\{u=v\}}+O(k^{-1}).
\]
En consecuencia $\iota_e^TS(s)^+\iota_{e'}$ es el número de vértices de fila y columna comunes a
$e,e'$, con
error $O(k^{-1})$.
\end{lema}
\begin{proof}
Como las masas de fila y de columna valen uno,
\[
 S_0=\begin{pmatrix}I&P\\P^T&I\end{pmatrix}.
\]
Denotemos por $\Pi_+,\Pi_-$ los proyectores ortogonales sobre las dos rectas generadas por
$(\mathbf1_R,\mathbf1_C)$ y $(\mathbf1_R,-\mathbf1_C)$, respectivamente. Sus entradas tienen módulo
$1/(2k)$. Los autovalores correspondientes son dos y cero. Sea $N_0$ la matriz simétrica de bloques
diagonales nulos y bloques cruzados $P-\mathbf1\mathbf1^T/k$ y su traspuesta. Anula ambas rectas
constantes, y en su complemento $S_0=I+N_0$. Tanto las filas como las columnas de $P$ tienen norma
euclídea $O(k^{-1/2})$, por $\sum_jp_{ij}^2\le C/k$ y su análoga en columnas. Vale lo mismo para las
filas de $N_0$. Sus entradas son $O(1/k)$ y $\|(I+N_0)^{-1}\|_{\rm op}\le1/\min\{1,\gamma\}$. La
identidad exacta
\[
 (I+N_0)^{-1}= I-N_0+N_0(I+N_0)^{-1}N_0
\]
muestra, por Cauchy--Schwarz entre dos filas de $N_0$, que sus entradas difieren de las de $I$ en
$O(k^{-1})$. Para recuperar $S_0^+$ se resta $\Pi_-+\tfrac12\Pi_+$, lo que respeta la misma
estimación.

Para $d=2,4$ se tiene $0\le\tau_d(p_e)-p_e=O(p_e^2)$. En cada fila de $S_{R,d}-S_0$ la suma absoluta
es dos veces la suma de estas diferencias sobre una línea física y es $O(k^{-1})$. Por simetría la
norma de operador tiene esta misma cota, también para $S(s)-S_0$. Todas las matrices tienen el mismo
núcleo y dominan $S_0$ en su complemento. En ese complemento la identidad de resolventes da
$\|S(s)^+-S_0^+\|_{\rm op}\le\gamma^{-2}\|S(s)-S_0\|_{\rm op}=O(k^{-1})$.
Al contraer con $\iota_e,\iota_{e'}$ intervienen sólo cuatro entradas. Esto da la afirmación restante.
\end{proof}
